\documentclass[11pt]{article}
\usepackage[margin=1in]{geometry}
\usepackage{times}
\usepackage{titlesec}
\usepackage{setspace}
\usepackage{parskip}
\usepackage{microtype}

\titleformat{\section}{\normalfont\bfseries\large}{\thesection}{1em}{}
\titlespacing*{\section}{0pt}{12pt}{6pt}

\begin{document}
\pagestyle{empty}

\begin{center}
{\LARGE\bfseries Constraint Optimization through RTL Enhancement Using Generative AI}\\[6pt]
{\large Team Timing Blinders --- Shubhanshu Shivhare, Nagarjun BV}\\[4pt]
{\small Astera Labs Nebula 2026, Digital Track}
\end{center}

\vspace{8pt}

\section*{Abstract}

A physical design engineer closing timing on a block spends most of the
cycle doing the same four things: read the STA report, find the path that
missed, pick a fix --- pipeline it, retime a register, restructure the logic,
or re-encode an FSM --- edit the RTL by hand, and re-synthesize to see if it
actually helped. None of that decision loop is inherently manual; it is
manual because nobody has trusted an automated system enough to run it
unsupervised. We build that system, and put the trust question first instead
of last.

Three tools, one loop. \texttt{run\_sta} wraps OpenSTA and turns a raw timing
report into ranked critical paths with slack, in a form an LLM can actually
read. \texttt{propose\_and\_apply\_edit} hands that path, and only that path's
RTL, to the model, constrained to emit exactly one of the four techniques
above as edited Verilog --- never prose, never a suggestion the model then
has to also implement correctly by hand. \texttt{resynth\_and\_check}
re-synthesizes the edit in Yosys, re-runs OpenSTA to confirm slack actually
improved, and then --- before the edit is allowed anywhere near floorplanning
--- checks it against the original RTL with EQY/SymbiYosys. Fail equivalence
and the edit is discarded, no exceptions, no bypass. Papers in this space
(ViTAD, SymRTLO, TimingLLM) get as far as diagnosing a violation and
proposing a fix; none of them formally prove the fix didn't break the
design before calling it a result. That gap is where this project sits.

The benchmark is sized to make the loop earn its numbers, not a toy that
happens to synthesize cleanly: five independent asynchronous master clock
domains, at least one generated clock per master, clock-domain crossings
between them, clock dividers at more than one ratio, roughly 50{,}000 standard
cells post-synthesis. It runs end to end through OpenROAD-flow-scripts on
sky130hd, so every accepted edit ships with matched before/after timing,
frequency, and PPA numbers alongside its equivalence-check trail --- not just
a slack delta.

Results get reported as measured, not as hoped for: how many of the injected
critical paths close inside a fixed iteration budget, how much slack comes
back on the ones that do, and which ones don't converge at all. ViTAD's
73.68\% repair rate against a 54.38\% plain-LLM baseline is the number we
compare against, honestly, rather than around.

\end{document}
